For odd $D\leq23$, the flag calculation uses one more restriction
on the simple support.

\begin{lemma}\label{lem:finite-k4}
Suppose that $D\leq23$ is odd and $m=f(D)+1$ marked edges form a
strong clique. No four vertices are joined pairwise by marked edges.
\end{lemma}

\begin{proof}

Suppose $D\le23$ is odd and $m=f(D)+1$. Let $A$ be the
vertex set of a copy of $K_4$ whose six pairs are marked. In particular,
$D\ge3$.
Complete the host to degree $D$ by adding unmarked edges
as in \textbf{Subsection~\ref{sec:weighted-counts}}. For each vertex $v$ put
$q_v=|N_G(v)\cap A|$ and $h_v=d_H(v)$. Let $J_A$ be
the sum, over the vertices of $A$, of the difference
between host degree and simple host degree. Let $m_A$
be the number of marked edges with both endpoints in
$A$. For $v\notin A$, let $r_v$ and $c_v$ count the
marked and unmarked edges from $v$ to $A$, respectively.
Then $r_v+c_v\ge q_v$ and $\sum_v q_v=4D-J_A$.

For a marked edge with both endpoints outside $A$, the
union of their neighborhoods in $A$ is a vertex cover of the
$K_4$ on $A$ and therefore has size at least three. For a
marked edge with exactly one endpoint in $A$, the sum
of $q_v$ over its endpoints is $3+q_v$, where $v$ is its
endpoint outside $A$. For a marked edge inside $A$ this
sum is six. Adding over all marked edges gives
\[
 3m+\sum_{a\in A}h_a+m_A+
 \sum_{v\notin A}(q_v-1)r_v\le\sum_v q_vh_v.
\]
Likewise, summing the endpoint values of $q_v$ over
unmarked edges gives
\[
 \sum_{a\in A}(D-h_a)+\sum_{v\notin A}(q_v-1)c_v
 \le\sum_v q_v(D-h_v).
\]
Here all edges inside $A$ are marked. An unmarked edge
from $A$ to $v\notin A$ contributes $3+q_v$ on the right
and $q_v$ on the left, and any unmarked edge with both
endpoints outside $A$ contributes only on the right.
Adding the two inequalities and using $m_A\ge6$ gives
\begin{equation}\label{eq:multi-k4-loss}
 3m\le4D(D-1)-6-DJ_A-\sum_{v\notin A}q_v(q_v-1).
\end{equation}

Put $W=\{v\notin A:q_v\ge2\}$. Every marked edge
avoiding $A$ meets $W$, and each vertex of $W$ has at
most $D-2$ incident edges avoiding $A$. There are at
most $4D-6$ marks meeting $A$. Thus
\[
 |W|\le\left\lfloor\frac{4D(D-1)-6-3m}{2}\right\rfloor,
 \qquad m\le4D-6+(D-2)|W|.
\]
The first bound is negative for odd $D\le11$. For odd
$13\le D\le21$, substituting the first bound into the
second gives $m\le m_{\rm up}$, with values
\[
\begin{array}{c|rrrrr}
 D&13&15&17&19&21\\ \hline
 m&206&275&354&443&542\\
 m_{\rm up}&46&106&212&342&534.
\end{array}
\]
In each column $m_{\rm up}<m$.

For $D=23$, let $n_i$ count the vertices outside $A$
with $q_v=i$, for $1\le i\le4$, and put
$g=65-2n_2-6n_3-12n_4$. By \eqref{eq:multi-k4-loss},
$g\ge0$. The difference $g$ bounds the number of marks
inside $W$, since each contributes at least four to the
sum of its endpoint values of $q_v$ and is charged only
three in the marked-edge inequality. The same difference
bounds the number of unmarked edges from vertices with
$q_v=0$ to vertices outside $A$ with $q_v>0$, since each
contributes at least one to the right of the unmarked-edge
inequality.

At least $651-(4\cdot23-6)=565$ marks avoid $A$.
Counting them at $W$ gives $21n_2+20n_3+19n_4\ge565$.
Every such mark either meets a vertex with $q_v=1$ or
$q_v\ge3$, or has both endpoints in $W$. Thus
$22n_1+20n_3+19n_4+g\ge565$. At most 20 distinct
neighbors outside $A$ meet each anchor, so
$n_1+2n_2+3n_3+4n_4\le80$. To list the possibilities, fix
nonnegative integers $n_2,n_3,n_4$ with $2n_2+6n_3+12n_4\le65$
and $21n_2+20n_3+19n_4\ge565$. The lower and upper bounds for
$n_1$ are then
\[
 \left\lceil\frac{565-20n_3-19n_4-g}{22}\right\rceil
 \le n_1\le80-2n_2-3n_3-4n_4 .
\]
The integer solutions have $n_3\in\{0,1,2\}$,
$n_1=26-n_3$, $n_2=27-n_3$, $n_4=0$ and $g=11-4n_3$.
The three cases are drawn in \textbf{Figure~\ref{fig:finite-k4-neighborhoods}}
of the supplement~\cite{FGSTSupplement}.
In each case every vertex of $A$ has 20 distinct neighbors outside
$A$, so every edge meeting $A$ is single.
The total capacity of $W$ for edges avoiding $A$ is
$567-n_3$. If $p$ counts marks within $W$, then
$p\le2-n_3$, because these marks are counted twice in
the capacity sum. There are at most $567-n_3$ marks
avoiding $A$, so at least $78+n_3$ of the 80 spokes at
$A$ are marked.

If $n_3$ is positive, there are at least $565-22n_1-p$
marks between vertices with $q_v=0$ and vertices with
$q_v=3$. This gives at least 14 marks when $n_3=1$ and
37 when $n_3=2$. As each vertex with $q_v=3$ has at
most 20 units of host degree outside $A$, it has at
most six or three units left for other exterior edges,
respectively. Choose one of these marked edges $xw$
with $q_x=0$ and $q_w=3$, and let $c$ be the unique vertex of
$A$ nonadjacent to $w$. For every marked spoke $cu$, its endpoint
$u$ must be adjacent to $x$ or $w$. Apart from $w$,
at most $n_3-1$ marked neighbors of $x$ have positive
$q_v$, and at most $g$ unmarked neighbors of $x$ have
positive $q_v$. The available degree at $w$ gives at
most six or three further neighbors with positive $q_v$.
Thus there are at most 13 or seven possible endpoints
of these spokes, respectively. But $c$ has at least
$18+n_3$ marked spokes, a contradiction.

If $n_3=0$, retain the host on the vertices with $q_v=1$
or $q_v=2$, keeping the marks only on edges between the
two sets. There are at least $565-p\ge563$ marked
edges. The sum of the endpoint degrees of each such edge in this
host is at most $22+21=43$.
On the side with $q_v=2$, the total degree available
outside $A$ is 567. At least 565 units are used to
count each exterior mark once. The $p$ marks internal
to the side use $p$ more units, and each unmarked
internal pair uses two units. Thus there are at most
$p+\lfloor\frac{2-p}{2}\rfloor\le2$ internal host pairs
on this side. On the other side the capacity is
$26\cdot22=572$ and at least 563 units are used for
crossing marks, leaving at most four internal host
pairs. We use the following bound for this bipartite marked
subgraph. If the Ore-degree of a bipartite marked graph is at most
$\sigma$ in its host, and at most $E$ pairs on each side are adjacent
in the host, then the number of marked edges is at most
$\frac{\sigma^2}{4}+\sigma\sqrt{E}$.

To prove the bound, take a vertex $v$ of maximum marked degree $a>0$,
of host degree $d$, and call its side $X$ and the other side $Y$.
Let $P$ be its $p$ distinct marked neighbors, $B_0$ the vertices
outside $N_G[v]$, and $H_0$ the marked graph on $B_0$.
Write $T$ for the vertices of $X\cap B_0$ incident with edges
of $H_0$. For each $u\in P$, put $T_u=T\setminus N_G(u)$ and
$Y_u=N_{H_0}(T_u)$. If $x\in T_u$ and $xy$ is marked in $H_0$,
then $uy$ must be an edge of $G$ to join $uv$ and $xy$. These are
unmarked edges on side $Y$. Hence $E'=\sum_{u\in P}|Y_u|\leq E$.

Choose weights $0\leq w_x\leq1$ for $x\in T$ satisfying
$d_{H_0}(x)\leq w_x(a-d_H(x,P))$. Put $W=\sum_{x\in T}w_x$
and $\delta=\sum_{u\in P}\sum_{x\in T_u}w_x-E'$.
Every vertex of $P$ has host degree at most $\sigma-d$.
Counting its neighbors in $T$ with weights $w_x$ and its neighbors
in $Y_u$ with weight one, and summing over $u$, gives
$pW-\delta\leq p(\sigma-d)$.
At most $a(d-a)$ marked edges meet $N_G(v)\setminus P$.
All remaining marked edges lie in $H_0$ or meet $P$.
For each adjacent pair $xu$ with $x\in T$ and $u\in P$,
either at least one unit of degree at $u$ is unmarked or
$w_x d_H(x,P)$ counts its weight in the bound for $d_{H_0}(x)$.
Thus, writing $m_0$ for the number of marks in the host, we obtain
\[
 m_0\leq a(d-a)+(a-p)W+p(\sigma-d)+\delta
 \leq a(\sigma-a)+\frac{a}{p}\delta.
\]
The second inequality uses $a\geq p$.

Fix $0<\theta<1$. Temporarily unmark each edge of $H_0$ whose
endpoint $x\in X$ satisfies $d_H(x,P)>(1-\theta)a$. There are at
most $\frac{p(\sigma-d)}{(1-\theta)a}$ such vertices, each incident
with at most $\theta a$ edges of $H_0$. For the remaining marks,
choose $w_x=\frac{d_{H_0}(x)}{a-d_H(x,P)}$ on the resulting $T$.
Then $w_x\leq\frac{d_{H_0}(x)}{\theta a}$. For each $y\in Y\cap B_0$,
let $g_y$ be the number of $u\in P$ such that $y\in Y_u$, with
$Y_u$ computed after unmarking. If $xy$ is still marked in $H_0$,
then every $u\in P\setminus N_G(x)$ contributes to $g_y$. As
$d_{H_0}(y)\leq a$, summing over marked edge copies gives
\[
 \sum_{u\in P}\sum_{x\in T_u}w_x
 \leq\frac{1}{\theta a}\sum_{xy\in E(H_0),\,x\in X}g_y
 \leq\frac{E'}{\theta}.
\]
Consequently $\delta\leq E(\frac{1}{\theta}-1)$. Adding back the
removed marks bounds their original total by
\[
 a(\sigma-a)+p(\sigma-d)\frac{\theta}{1-\theta}
              +\frac{aE}{p}\frac{1-\theta}{\theta}.
\]
Taking the infimum over $\frac{\theta}{1-\theta}>0$ gives
$a(\sigma-a)+2\sqrt{a(\sigma-d)E}$. Since $a\le d$ and
$a(\sigma-a)\le\frac{\sigma^2}{4}$, this proves the stated bound.
For $\sigma=43$ and $E=4$, it gives at most
$\frac{43^2}{4}+86=\frac{2193}{4}<563$ marked edges. This
contradicts the required number of crossing marks.
\end{proof}
